\documentclass[12pt,a4paper]{article}

\usepackage[a4paper,margin=1in]{geometry}
\usepackage{hyperref}
\usepackage{enumitem}

\title{Computer Science Portfolio Project Report}
\author{Arjun Bijoy \ Student ID: GH1044696}
\date{Submission Deadline: June 26, 2026}

\begin{document}

\maketitle

\begin{center}
GitHub Repository Link: \
\url{https://github.com/arjunbijoy1917r-cmyk/computer-science-lab.git}

\vspace{0.3cm}

Portfolio Website Link: \
\url{https://arjunbijoy1917r-cmyk.github.io/computer-science-lab/portfolio_site/}
\end{center}

\newpage

\tableofcontents

\newpage

\section{Introduction}

This report explains the design, implementation, and content of my computer science portfolio website. The main aim of this portfolio is to show my skills, learning progress, and projects as a Software Engineering student.
The portfolio website was developed using HTML and CSS. It includes personal information, achievements, completed work, technical skills, and contact information. The website is also connected to my GitHub repository so visitors can view my project files and source code.
The portfolio is hosted using GitHub Pages. This makes the website available online through a public URL. Overall, this report explains how the portfolio was built, why the design choices were made, and what improvements can be added in the future.

\section{Website Structure}

The portfolio website has a simple structure. The navigation bar at the top contains links to the main sections of the page. These sections are Home, About, Achievements, Work, Skills, and Contact Me.
This structure was selected because it makes the website easy to understand and easy to browse.

\subsection{Home Section}

The Home section is the first section of the portfolio. It includes a greeting, my name, and a short introduction about me. It gives visitors a quick idea about who I am and what I am studying.

\subsection{About Section}

The About section explains my background as a Software Engineering student. It also mentions my interest in programming, web development, and databases.
This section shows that I am currently learning and improving my skills in HTML, CSS, Python, Java, SQL, and LaTeX.

\subsection{Achievements Section}

The Achievements section shows the progress I have made so far. It includes completed Python projects, SQL database projects, technical exercises, and beginner-level projects.
This section is useful because it shows my learning journey, not only the final result.

\subsection{Work Section}

The Work section presents my main projects. The two main projects included in my portfolio are:

\begin{itemize}
\item \textbf{Python University System:} A beginner-level Python project based on a university system. It uses concepts such as classes, objects, students, courses, instructors, and enrollments.


\item \textbf{University Database Project:} A SQL database project based on a university management system. It includes tables, sample data, queries, relationships, and database design.


\end{itemize}

Each project has a button that opens the related GitHub repository. This makes it easy for visitors to explore the project work.

\subsection{Skills Section}

The Skills section shows the main technologies I am currently learning. These include:

\begin{itemize}
\item Basic Python Coding
\item Java
\item HTML
\item CSS
\item SQL Database
\item LaTeX
\end{itemize}

I kept this section short so visitors can quickly understand my current technical skills.

\subsection{Contact Section}

The Contact Me section includes my email address. This helps visitors contact me for questions or updates. The contact link is useful because it allows people to send an email directly.

\section{GitHub Repository Structure}

The GitHub repository is named \textbf{computer-science-lab}. Inside the repository, the portfolio website is stored in the \textbf{portfolio_site} folder.

The main files used in the portfolio website are:

\begin{itemize}
\item \textbf{index.html:} This file contains the main structure and content of the website.

```
\item \textbf{styles.css:} This file controls the design of the website, including colours, layout, spacing, borders, and hover effects.

\item \textbf{assets folder:} This folder stores images and other media files used in the website.

\item \textbf{projects.html:} This file can be used to show project-related content or can be improved later as a separate project page.
```

\end{itemize}

This structure is useful because it keeps the website files organised. It also makes it easier to edit the website. For example, if I want to change the design, I can edit the CSS file. If I want to change the content, I can edit the HTML file.
\section{Design Decisions}

The design of the portfolio was kept simple and readable. I did not make the website too complicated because the main purpose is to present my work clearly.

\subsection{Simple Navigation}
A navigation bar was added at the top of the website. This allows visitors to move easily between the main sections. The section names are simple and easy to understand.

\subsection{Colour Choice}

The website uses a light background with white sections and a brownish-orange accent colour. This colour is used for headings, borders, and buttons.
The light background improves readability, while the accent colour gives the website a personal and professional style.

\subsection{Section-Based Layout}

Each part of the portfolio is placed in a separate section. This makes the website easier to read because the information is not mixed together.
The section-based layout also gives the website a clean structure.
\subsection{Project Cards and Buttons}

The projects are displayed using project boxes. Each project includes a title, short description, and a button linking to GitHub.
This makes the portfolio more interactive and helps visitors open the project repositories directly.

\subsection{Beginner-Friendly Design}

The website was designed using basic HTML and CSS. This was done because it reflects my current learning stage.
Using simple code helped me understand each part of the website properly. The portfolio is not only a final product, but also part of my learning process.
\section{Tools and Technologies Used}

The following tools and technologies were used to create the portfolio:

\begin{itemize}
\item \textbf{HTML:} Used to create the structure and content of the webpage.


\item \textbf{CSS:} Used to style the website, including colours, fonts, spacing, layout, borders, and hover effects.

\item \textbf{GitHub:} Used to store the project files and manage the repository.

\item \textbf{GitHub Pages:} Used to host the portfolio website online.

\item \textbf{LaTeX:} Used to prepare this project report in a structured format.

\item \textbf{Web Browser:} Used to test the portfolio website after uploading it to GitHub Pages.


\end{itemize}

These tools were selected because they are suitable for a beginner-level computer science portfolio.

\section{Implementation Explanation}

The implementation started by creating the basic HTML structure. I divided the webpage into different sections so each part had a clear purpose.
After creating the HTML file, I used CSS to improve the appearance of the website. The CSS file controls the navigation bar, hero section, image placement, section cards, project boxes, buttons, and contact links.
Using a separate CSS file was a good decision because it keeps the design separate from the HTML content.After completing the website files, I uploaded them to the GitHub repository inside the portfolio folder. Then I used GitHub Pages to publish the website online.This process helped me understand how websites can be hosted using GitHub and how folder paths affect the final website link.

\section{Portfolio Content}

The content of the portfolio is based on my current learning and projects. The About section explains my background as a Software Engineering student. The Achievements section shows my progress, and the Work section shows my Python and SQL projects.
The Python University System project shows my understanding of basic programming and object-oriented concepts. The University Database Project shows my understanding of database tables, relationships, queries, and database design.
These projects are suitable for my portfolio because they are connected to computer science and software engineering.

\section{Strengths of the Portfolio}

One strength of the portfolio is that it has a simple and clear design. The sections are clearly labelled, and the navigation menu helps users move around the page easily.
Another strength is that the project buttons connect directly to GitHub. This allows visitors to check the actual project work.
The design is also consistent because the same accent colour is used throughout the website. The use of separate HTML and CSS files also makes the website easier to edit and improve.

\section{Weaknesses of the Portfolio}

The portfolio also has some weaknesses. The design is currently basic and does not include advanced features such as animations, improved responsive design, or a working contact form.
Some project descriptions can also be explained in more detail. At the moment, the website mostly includes text and a profile image. Adding screenshots, diagrams, or short explanations of the development process could make the portfolio more professional.

\section{Future Improvements}

In the future, I would like to improve the portfolio by adding more features. Some planned improvements are:

\begin{itemize}
\item Adding screenshots of each project.
\item Adding a separate project details page for each project.
\item Improving the mobile view for phones and tablets.
\item Adding a working contact form.
\item Adding my CV as a downloadable PDF.
\end{itemize}

These improvements can make the portfolio more professional and interactive. They can also show my progress as I gain more knowledge and technical skills.

\section{Conclusion}

In conclusion, this portfolio project helped me practise basic web development and learn how to publish a website using GitHub Pages.
The website presents my background, achievements, projects, skills, and contact information. The GitHub repository also shows the structure of the website and allows the source code to be reviewed.
Although the portfolio has a simple design, it is a good starting point for showing my computer science learning progress. The main strengths are its clear structure, simple navigation, project links, and consistent design.
The main areas for improvement are responsive design and additional interactive features. Overall, this project helped me understand the connection between website design, content organisation, and online hosting.

\end{document}

